\documentclass[11pt]{article}
\usepackage[margin=1in]{geometry}
\usepackage{amsmath,amssymb}
\usepackage[hidelinks]{hyperref}

\begin{document}

\section*{Human Review: convex-block homogeneous misreading}

\noindent Packet identifier:
\texttt{1711.01340\_ell2\_difference\_basis\_not\_submultiplicative}

\noindent Original automatic proof: \texttt{solution\_packet.pdf}

\noindent Checked date: 14 June 2026

\noindent Status: Manually checked.

\noindent Verdict: Rejected as a solution to the intended problem.  This is a
miss.

\section*{Human Comments}

The proposed solution in \texttt{solution\_packet.pdf} should not be accepted
as a counterexample to Problem 7 of Motakis--Puglisi--Tolias.  The finite
calculation with the difference sequence in \(\ell_2\) is correct, but the
example does not satisfy the hypothesis of the problem.

The packet claims that the usual unit vector basis of \(\ell_2\) is convex
block homogeneous because every normalized block sequence is orthonormal.
This uses the wrong notion.  In the context cited by the source paper,
Argyros--Motakis--Sari introduce a convex block homogeneous sequence as one
which is equivalent to its convex block sequences.  The convex blocks are not
renormalized.

For the unit vector basis of \(\ell_2\), take successive finite sets \(B_k\)
with \(|B_k|\to\infty\), and define
\[
  y_k=\frac{1}{|B_k|}\sum_{i\in B_k} e_i .
\]
Then \((y_k)\) is a convex block sequence of the unit vector basis, but
\[
  \|y_k\|_2=|B_k|^{-1/2}\longrightarrow 0.
\]
Therefore \((y_k)\) is not even seminormalized and cannot be equivalent to the
normalized unit vector basis of \(\ell_2\).  Thus the proposed example fails
the defining hypothesis ``convex block homogeneous basis''.

This also fits the surrounding context of the source paper.  Remark 8.10
mentions the known convex block homogeneous examples arising from \(J(X)\) and
\(J^*(X)\), and then says that a positive answer to Problem 7 would imply a
statement about difference bases of conditional spreading sequences.  The
question is not aimed at the Hilbert unit vector basis under a normalized-block
interpretation.

There is a secondary warning as well: in \(\ell_2\), the sequence
\[
  d_1=e_1,\qquad d_{i+1}=e_{i+1}-e_i
\]
is not a Schauder basis of \(\ell_2\).  The inverse coordinate map involves
partial summation and is unbounded.  This reinforces that the packet is testing
a coefficient norm outside the intended setting.

\section*{Short Version}

This is a miss.  The packet treats ``convex block homogeneous'' as equivalence
to normalized block sequences, but the source context uses equivalence to
convex block sequences themselves.  The \(\ell_2\) unit vector basis fails that
hypothesis by taking long equal-average convex blocks.

\section*{Review Outcome}

Reject this candidate as a proposed counterexample to Problem 7.  It may be
archived as a warning that automatic extraction must preserve the precise
definition of ``convex block homogeneous'' and the surrounding conditional
spreading-sequence context.

\end{document}
